\par\addvspace{12pt}\Needspace{9\baselineskip}

\section{随机过程初步*}

% ITEM: probability-example-036
\solutionitem{probability-example-036}{例6.36 *}
本节投篮模型：甲、乙命中率为\(0.6,0.8\)，命中者继续，未中则换人；每次命中率不受以往情况影响，首次由谁投篮各以\(1/2\)决定。

（据2023年新课标I卷改编）
沿用本节开头的投篮模型。

（1）求第二次由乙投篮的概率。

（2）求第\(i\)次由甲投篮的概率，并研究其极限。

（3）求前\(n\)次投篮中甲投篮次数的期望。说明这里为什么不能直接使用同一参数的二项分布。
\solutionanswer
（1）第二次由乙的概率\((1/2)(0.4)+(1/2)(0.8)=0.6\)。
（2）按正文索引，\(u_0=1/2\)对应第一次投篮；\(u_{n+1}=0.2+0.4u_n\)，固定点\(1/3\)。所以第\(i\)次由甲的概率为\(u_{i-1}=1/3+(1/6)(2/5)^{i-1}\)，极限\(1/3\)。矩阵法也可得到同式：\(\pi=(1/3,2/3)^{\mathsf T}\)满足\(M\pi=\pi\)，\(w=(1,-1)^{\mathsf T}\)满足\(Mw=(2/5)w\)。将初始分布写成\(\pi+w/6\)，再作用\(M^{i-1}\)即可。
（3）设每次由甲投篮的指示量为\(I_i\)，则次数期望
\[
 \sum_{i=1}^nEI_i=\frac n3+\frac5{18}\left[1-\left(\frac25\right)^n\right].
\]
投篮者受前一状态及命中结果影响，各次指示量并非独立，边际概率也随次数改变，不能直接作为同一参数的二项分布；期望线性性仍可使用。


% ITEM: probability-exercise-29-01
\solutionitem{probability-exercise-29-01}{习题6.29（1） *}
本组\(S_j\in\{A,B\}\)，\(u_j=P(S_j=A)\)，转移概率\(A\to B\)为\(\alpha\)、\(B\to A\)为\(\beta\)，且\(u_{j+1}=\beta+(1-\alpha-\beta)u_j\)。前\(n\)步计入\(j=0,\ldots,n-1\)。
取\(\alpha=0.3,\beta=0.2\)，初始状态为\(A\)。求\(u_n\)、不变分布及前\(n\)步处于\(A\)的次数期望。
\solutionanswer
\(u_*=0.2/(0.3+0.2)=2/5\)，\(r=1/2,u_0=1\)，所以\(u_n=2/5+(3/5)2^{-n}\)。唯一不变分布为\((2/5,3/5)^{\mathsf T}\)。按正文前\(n\)步指\(j=0,\ldots,n-1\)，次数期望\(2n/5+(6/5)(1-2^{-n})\)。


% ITEM: probability-exercise-29-02
\solutionitem{probability-exercise-29-02}{习题6.29（2） *}
本组\(S_j\in\{A,B\}\)，\(u_j=P(S_j=A)\)，转移概率\(A\to B\)为\(\alpha\)、\(B\to A\)为\(\beta\)，且\(u_{j+1}=\beta+(1-\alpha-\beta)u_j\)。前\(n\)步计入\(j=0,\ldots,n-1\)。
分别取\(\alpha=\beta=0\)与\(\alpha=\beta=1\)，比较不变分布的唯一性与分布收敛性。再说明\(\alpha=0,0<\beta<1\)时模型的长期行为。
\solutionanswer
\(\alpha=\beta=0\)时转移矩阵为恒等矩阵，每个概率向量都是不变分布，初始分布始终保持，故虽各自收敛，却没有统一且唯一的长期分布。\(\alpha=\beta=1\)时唯一不变分布为\((1/2,1/2)^{\mathsf T}\)，但\(u_n=1/2+(-1)^n(u_0-1/2)\)，除初始恰为该分布外不收敛。不变说的是分布不变，并非状态不再变动。
\(\alpha=0,0<\beta<1\)时\(A\)吸收，\(u_n=1-(1-u_0)(1-\beta)^n\to1\)，唯一不变分布为\((1,0)^{\mathsf T}\)，任意初始分布都趋向它。
\par
